\label{chap:Methodology}
This chapter describes the evaluation methodology used to benchmark the 
automated pipeline against the manually reviewed OGS reference catalog. The 
central challenge is to establish rigorous, one-to-one correspondences between 
automated detections and catalog entries, so that each detection is matched to 
at most one catalog pick or event. To this end, we employ a bipartite graph 
matching framework grounded in weighted similarity metrics.

In this formulation, detected picks and catalog entries form two disjoint node 
sets. Edge weights encode multi-criteria similarity scores that combine 
temporal proximity, phase consistency, station coverage, and physical 
plausibility. A maximum-weight matching algorithm then selects the globally 
optimal, non-conflicting set of correspondences, eliminating duplicate 
assignments and ensuring consistent benchmarking against the manually curated 
OGS catalog.


\section{Distance and Similarity Functions}

The matching framework relies on two core similarity functions, one temporal 
and one spatial, whose outputs are combined into a single weighted score.

\subsection{Time-Based Similarity}
The time similarity function \texttt{dist\_time()} computes a score based on 
the temporal difference between predicted and reference arrivals, normalized by 
the pick-weight uncertainty:

\begin{equation}
  S_t = \max\left(0, 1 - \frac{|t_{\text{pred}} - t_{\text{ref}}|}{\sigma_w}\right)
\end{equation}

where $\sigma_w$ is the uncertainty corresponding to the reference pick weight. 
A perfect temporal match yields $S_t = 1$, while differences exceeding 
$\sigma_w$ yield $S_t = 0$.

\subsection{Spatial Distance}
The spatial distance function \texttt{dist\_space()} computes geodetic 
distances using ObsPy's \texttt{gps2dist\_azimuth()} function, which implements 
the WGS84 ellipsoid model for accurate great-circle distance calculations.

\subsection{Weighted Matching Scores}
The overall matching score for each candidate pair is a weighted combination 
of temporal and spatial components. The relative weights are configurable, 
allowing the scoring scheme to be tuned for different operational priorities 
(see Sections~\ref{subsec:matching_criteria} for the specific weight values 
used in this study).

\section{Confusion Matrix}
Confusion matrices are a standard tool in \ac{ML} for evaluating classification 
performance. Table~\ref{tab:cfn_mtx} shows the confusion matrices used to 
evaluate the seismic phase pickers and associators used in this 
study. A confusion matrix is a tabular representation that summarizes the 
performance of a classification model by comparing two classes of observations. 
Each entry $C_{ij}$ records the number of instances assigned to predicted 
\textbf{Target} class ($T_j$), given that the corresponding \textbf{Base} class 
($B_i$) is known. In the seismic context, the $T_j$ are the outputs of the 
automated phase pickers (\textit{e.g.} PhaseNet and EQTransformer) and 
associators (\textit{e.g.} GaMMA and PyOcto), while the $B_i$ are the P- and 
S-wave picks (and their associated events) in a hand-curated catalog (here, the 
OGS reference catalog). Thus, the confusion matrix quantifies how many catalog 
picks or events are correctly recovered by the automated method versus how many 
are missed, and how many are spuriously reported.

In applying confusion matrices to seismic pick detection, analogous to the 
\ac{ML} domain in confusion matrices for classification, we propose a new 
standard terminology to reflect the seismic phase picking and event association 
context. For example, a “\textbf{True Positive}” in ML simply means a correct 
classification, but in seismic picking we use the term “\textbf{Matched}” (MH) 
to emphasize a specific, criterion-based (see \ref{subsec:matching_criteria}) 
correspondence between a $B_i$ and a $T_j$. Mathematically, this can be 
expressed as:
$$B_i \Leftrightarrow T_j$$
Similarly, what ML calls a "\textbf{False Negative}" (an instance where the 
model missed a positive case) we term a “\textbf{Missed}” (MS) pick or event: 
these are catalog entries that our method failed to detect or failed to match 
under our strict criteria. "\textbf{False Positives}" are renamed to 
"\textbf{Proposed}" (PS) picks or events, indicating those that were detected 
by the automated method but are not present in the OGS reference catalog, which 
could be due to false detections or very subtle microseismic events not 
recorded in the manually reviewed catalog. The terms $\epsilon_{PS}$, and 
$\epsilon_{SP}$, represents proposed S picks correspond to P picks in the 
catalog, and vice versa, due to phase misclassification errors. The 
$\epsilon_{NN}$ term is provided for completeness, and can be interpreted as 
the number of instances where both the catalog and the automated method agree 
that no pick/event is present. However, since the absence of a pick in a time 
continuum framework is not explicitly cataloged, this term is not directly 
measurable and is often omitted from performance metrics. By adopting this 
terminology, we can more intuitively interpret the confusion matrices in the 
context of seismic phase picking and event association, while still leveraging 
the powerful analytical framework provided by \ac{ML} evaluation metrics.\\

At the pick level, the confusion matrices report Matched (MH), Missed (MS), and 
Proposed (PS) counts for P and S waves across all model--dataset combinations 
(INSTANCE, STEAD, SCEDC, and each model's original training set). Diagonal 
cells represent correctly identified picks~(MH); off-diagonal cells capture 
misclassifications and missed detections. Color intensity (logarithmic scale) 
is proportional to the relative frequency of each outcome, enabling rapid 
visual comparison of model behavior. At the event level, the confusion matrices 
extend this assessment to the full detection pipeline (picking + association), 
since successful event detection requires multiple coherent phase picks across 
the seismic network.

\begin{table}[h!]
  \centering
  \begin{tabular}{cc||c|c|c|}
    \hline
    \multirow{3}{*}{\rotatebox{90}{\textbf{Base}}} & P    & $MH_{P}$        & $\epsilon_{PS}$ & $MS_P$       \\\cline{2-5}
    & S    & $\epsilon_{SP}$ & $MH_{S}$        & $MS_S$       \\\cline{2-5}
    & None & $PS_P$          & $PS_S$          & $\epsilon_{NN}$ \\\hline\hline
    &      & P               & S               & None \\
    &      & \multicolumn{3}{c|}{\textbf{ML (Target)}} \\
  \end{tabular} $\qquad$
  \begin{tabular}{cc||c|c|}
    \hline
    \multirow{2}{*}{\rotatebox{90}{\textbf{Base}}}  & Event & $MH_{E}$ & $MS_{E}$        \\\cline{2-4}
                                                    & None  & $PS_{E}$ & $\epsilon_{NN}$ \\\hline\hline
                                                    &       & Event & None \\
                                                    &       & \multicolumn{2}{c|}{\textbf{ML (Target)}} \\

  \end{tabular}
  \caption{Confusion matrices for Picks (left) and Events (right).
  Manual OGS catalog (Base); automated \ac{ML} output (Target). MH: Matched; 
  MS: Missed; PS: Proposed.}
  \label{tab:cfn_mtx}
\end{table}

\subsection{Matching Criteria for Picks and Events}
\label{subsec:matching_criteria}
Rigorous evaluation of both phase picking and event association requires a 
transparent and reproducible protocol for determining when an automated output 
corresponds to a catalog entry. To ensure a physically meaningful assessment of 
system performance, we define an explicit set of criteria for what constitutes 
a ``\textbf{Match}'' between a detected pick/event and a catalog entry. These 
criteria balance \textbf{strictness} (to avoid overestimating performance) and 
\textbf{tolerance} (to avoid penalizing near misses due to observational 
uncertainty). The matching criteria below impose temporal and spatial 
tolerances grounded in the physics of seismic wave propagation and in the 
practical uncertainties of automated processing.

With these criteria, each detection is classified as Matched~(MH), Missed~(MS), 
or Proposed~(PS) based on its correspondence to the catalog. Ambiguities often 
arise when multiple detections fall within close proximity to a single catalog 
entry. To resolve these ambiguities, we adopt a bipartite matching scheme that 
selects the pairing that maximizes the composite matching score 
($\Delta_{\text{Picks}}$ or $\Delta_{\text{Events}}$). The matching score is a 
weighted sum of measurables such as temporal offset~($\delta t$), phase 
consistency~($\delta p$), pick certainty~($\delta \rho$), and, for events, 
spatial discrepancy~($\delta d$). We adjust each measurable's influence in the 
overall matching score by multiplying  them with predefined 
coefficient~($w_i$). These weights can be tuned to emphasize the sensitivity of 
the matching procedure to different aspects of pick or event quality, depending 
on the operational priorities or scientific objectives.

\subsubsection{Pick Matching}
A detected P- or S-wave pick is classified as a "Match" (MH) to a catalog pick 
if it satisfies all of the following conditions:
\begin{enumerate}
  \item \textbf{Station Consistency:} The detected and catalog pick originate 
  from the same seismic station, ensuring that the comparison is evaluated 
  based on each station's location and instrumentation.
  \item \textbf{Phase Agreement:} The detected phase type matches the catalog 
  phase (P or S), preserving physical interpretability in arrival time 
  comparisons and ensuring that P-wave picks are not erroneously matched to 
  S-wave picks, and vice versa.
  \item \textbf{Temporal Proximity:} The detected pick time falls within a time 
  window $\omega_t = 0.5$ seconds of the catalog pick time. This time interval 
  accounts for typical picking uncertainties from noise and variable signal 
  quality.
\end{enumerate}

Additionally, we incorporate confidence information when available. Both 
automated pickers and OGS analysts provide a pick certainty parameter 
($\delta \rho$), enabling a broader measure of pick-level similarity. We define 
a composite score:
$$\Delta_{Picks} = \sum_{i\in \left\{t, p, \rho\right\}} w_i \delta_i$$
where
\begin{itemize}
  \item $\delta t$ represents the absolute arrival time difference;
  $$\delta t = \begin{cases}
    1 - \frac{|T_t - B_t|}{\omega_t}, & \text{if } |T_t - B_t| \leq \omega_t \\
    0, & \text{otherwise}
  \end{cases}$$
  where $T_t$ and $B_t$ are the detected and catalog pick times, respectively, 
  and $\omega_t$ is the predefined temporal tolerance (0.5 seconds).
  \item $\delta p$ is a binary variable indicator quantifying whether both 
  picks correspond to the same phase type (P or S).
  $$\delta p = \begin{cases}
    1, & B_p = T_p\\
    0, & B_p \neq T_p
  \end{cases}$$
  where $T_p$ and $B_p$ are the detected and catalog pick phase types,
  respectively. Using a binary indicator for phase type mismatch ensures that 
  picks of different phases are heavily penalized in the overall matching 
  score.
  \item $\delta_\rho$ compares the confidence values assigned to the 
  detected and catalog picks, incorporating probabilistic information about 
  pick reliability (\textit{e.g.} signal clarity, analyst confidence).
  $$\delta \rho = \frac{T_\rho}{B_\rho}$$
  where $T_\rho$ and $B_\rho$ are the confidence values assigned to the 
  detected and catalog picks, respectively.
  \item $w_i$ are adjustable non-negative coefficients that control the 
  relative influence of each component, satisfying the normalization 
  condition $\sum_i w_i = 1$. In this study:
  $$w_t=0.97, \quad w_p=0.02, \quad w_\rho=0.01$$
  The dominant weight on $\delta t$ reflects the primacy of temporal accuracy 
  in phase-pick evaluation.
\end{itemize}

\subsubsection{Event Matching}
An automatically detected seismic event is a coherent set of associated picks 
spanning multiple stations. A detected event~($T_j$) is classified as a 
``Match''~(MH) to a catalog event~($B_i$) if it satisfies the following 
spatiotemporal criteria:
\begin{enumerate}
  \item \textbf{Temporal Proximity}: The origin time discrepancy between the 
  detected and catalog events lies within a tolerance of $\omega_t = 2$ 
  seconds. This threshold accommodates uncertainties arising from phase picking 
  variability, network coverage, and event location inversion.
  \item \textbf{Spatial Proximity}: The epicentral distance between the 
  detected event and the catalog event must be less than 8 km, $\omega_d = 8$ 
  km. The distance is calculated as the Euclidean epicentral distance between 
  the detected event location ($T_d$) and the catalog event location ($B_d$) 
  under projected coordinates. This constraint reflects typical lateral 
  uncertainties for regional networks with moderate station density and 
  acknowledges the variability inherent in automated location algorithms.
\end{enumerate}
These tolerances are consistent with expected uncertainties arising from noise, 
inversion errors, and variability in modern automated location algorithms. To 
quantify the similarity between detected and catalog events, we define the 
overall event-level similarity score as:
$$\Delta_{Events} = \sum_{i\in \left\{t, d\right\}} w_i \delta_i$$
where
\begin{itemize}
  \item $\delta t$ is the absolute origin time difference;
  $$\delta t = \begin{cases}
    1 - \frac{|T_t - B_t|}{\omega_t}, & \text{if } |T_t - B_t| \leq \omega_t \\
    0, & \text{otherwise}
  \end{cases}$$
  where $T_t$ and $B_t$ are the predicted and catalog event times, 
  respectively, and $\omega_t$ is the predefined temporal tolerance (2
  seconds).
  \item $\delta d$ is the epicentral distance between detected and catalog 
  event locations, calculated using Euclidean distance in Mercator projected
  coordinates.
  $$\delta d = \begin{cases}
    1 - \frac{||T_d - B_d||_2}{\omega_d},&\text{if }||T_d-B_d||_2\leq\omega_d\\
    0, & \text{otherwise}
  \end{cases}$$
  where $T_d$ and $B_d$ are the predicted and catalog event locations, 
  respectively, and $\omega_d$ is the predefined spatial tolerance (8 km).
  \item $w_i$ are the adjustable weights assigned to each parameter to 
  balance their influence on the overall matching score.
  $$\sum_i w_i = 1$$
  $$w_t=0.99, \quad w_d=0.01$$
\end{itemize}

Together, these criteria provide a structured, physically informed framework 
for establishing one-to-one correspondences between detected and cataloged 
phases and events, while accommodating the uncertainties inherent in real-world 
seismic data processing.

\section{Time Residual Analysis}
For each matched pick or event, the time residual is defined as 
$\Delta t = T_t - B_t$, i.e., the difference between the detected and catalog 
arrival (or origin) times. A narrow, symmetric, zero-centered residual 
distribution indicates high timing fidelity and minimal systematic bias; a 
broader or skewed distribution signals systematic timing errors or increased 
uncertainty.

Two complementary error metrics are reported: the Mean Absolute Error~(MAE), 
which is robust to outliers, and the Root Mean Square Error~(RMSE), which is 
sensitive to large deviations~(\cite{WhichPicker}). For each model 
configuration, we also report the mean~($\mu$) and standard 
deviation~($\sigma$) of the residual distribution to quantify bias and spread, 
enabling direct comparison across model and parameter choices.

\section{Event Location Accuracy}
For each matched event, we compute the epicentral distance between the 
detected and reference (OGS) locations. This metric captures spatial 
discrepancies arising from station geometry, velocity-model assumptions, and 
phase-arrival uncertainties. Combined with time residuals, spatial residuals 
reveal systematic propagation errors that can guide velocity-model refinement 
and network optimization. Mapping spatial residuals as a function of event 
depth or azimuthal gap further exposes network-dependent biases in the 
automatic locator.

\section{Magnitude Residual Analysis}
For each matched event, the magnitude residual is defined as:
\begin{equation*}
  \Delta M_L = M_{L : \text{detected}} - M_{L : \text{catalog}}
\end{equation*}
where $M_{L : \text{detected}}$ is the local magnitude estimated by our
automated pipeline, and $M_{L : \text{catalog}}$ is the local magnitude 
reported in the OGS catalog.

Positive residuals ($\Delta M_L > 0$) indicate overestimation; negative 
residuals indicate underestimation. We examine the distribution of magnitude 
residuals to assess systematic biases (mean) and random scatter (standard 
deviation). Additionally, we analyze how residual statistics and the balance 
between magnitude recall and \ac{FDR} vary with the detection threshold, and 
we report Missed and Proposed event counts as a function of threshold to 
characterize the sensitivity--specificity trade-off.

\section{Machine Learning Metrics}
In order to provide a comprehensive evaluation of our seismic phase picking and 
event association algorithms, we employ several key performance metrics 
commonly used in machine learning, adapted to the seismological context. We 
classify the outcomes of our detection algorithms into three categories based 
on the confusion matrix terminology introduced earlier: Matched (MH), Missed 
(MS), and Proposed (PS). We then derive the following metrics to quantify 
performance at both the pick and event levels:
\begin{itemize}
  \item \textbf{Recall}: 
    $$\mathrm{Recall}_P = \dfrac{\mathrm{MH}_P}{\mathrm{MH}_P + \mathrm{MS}_P + \epsilon_{PS}}$$
    $$\mathrm{Recall}_S = \dfrac{\mathrm{MH}_S}{\mathrm{MH}_S + \mathrm{MS}_S + \epsilon_{SP}}$$
    $$\mathrm{Recall}_{\text{Picks}} = \dfrac{\mathrm{MH}_P + \mathrm{MH}_S}{\mathrm{MH}_P + \mathrm{MH}_S + \mathrm{MS}_P + \mathrm{MS}_S + \epsilon_{PS} + \epsilon_{SP}}$$
    $$\mathrm{Recall}_{\text{Events}} = \dfrac{\mathrm{MH}_E}{\mathrm{MH}_E + \mathrm{MS}_E}$$
    Recall measures the fraction of catalog entries correctly recovered by 
    the automated method. At pick level, it quantifies the fraction of catalog 
    picks recovered; at event level, the fraction of catalog earthquakes 
    detected. High recall is essential to avoid overlooking significant events, 
    particularly in early-warning and hazard-monitoring contexts 
    (\cite{SeismicPhaseDetection}).
  \item \textbf{False Detection Rate (\ac{FDR})}:
    $$\mathrm{\ac{FDR}}_P = \dfrac{\mathrm{PS}_P}{\mathrm{MH}_P + \mathrm{PS}_P + \epsilon_{SP}}$$
    $$\mathrm{\ac{FDR}}_S = \dfrac{\mathrm{PS}_S}{\mathrm{MH}_S + \mathrm{PS}_S + \epsilon_{PS}}$$
    $$\mathrm{\ac{FDR}}_{\text{Picks}} = \dfrac{\mathrm{PS}_P + \mathrm{PS}_S}{\mathrm{MH}_P + \mathrm{MH}_S + \mathrm{PS}_P + \mathrm{PS}_S + \epsilon_{PS} + \epsilon_{SP}}$$
    $$\mathrm{\ac{FDR}}_{\text{Events}} = \dfrac{\mathrm{PS}_E}{\mathrm{MH}_E + \mathrm{PS}_E}$$
    \ac{FDR} measures the proportion of automated detections absent from the 
    reference catalog. Low \ac{FDR} is essential to avoid alarm fatigue, 
    preserve analyst trust, and ensure resource-efficient response. Although 
    Precision and F1 are standard metrics in the \ac{ML} literature, this 
    work focuses on recall, \ac{FDR}, and timing residuals as the primary 
    evaluation criteria.
  \item \textbf{Timing residuals}: Mean ($\mu$), standard deviation ($\sigma$), 
    Root Mean Square Error (RMSE) and Mean Absolute Error (MAE) of pick and 
    event origin-time residuals. Lower spread improves association and location 
    stability.
\end{itemize}

\section{Bipartite Graph Matching}

The matching scores defined above are used within a bipartite graph 
maximum-weight matching algorithm, implemented with the NetworkX library.

\subsection{Graph Construction}
The algorithm constructs a bipartite graph in which:
\begin{itemize}
  \item Nodes in partition~$A$ represent reference (catalog) picks or events
  \item Nodes in partition~$B$ represent predicted (automated) picks or events
  \item Edge weights encode the composite similarity score between each pair
\end{itemize}

\subsection{Matching Classes}
Two specialized classes implement the matching:
\begin{itemize}
  \item \textbf{OGSBPGraphPicks}: Maximum-weight matching for phase picks, 
  enforcing station and time constraints.
  \item \textbf{OGSBPGraphEvents}: Maximum-weight matching for seismic 
  events, incorporating spatiotemporal similarity and magnitude differences.
\end{itemize}
Both classes use the Hungarian algorithm to guarantee an optimal one-to-one 
correspondence between predicted and reference observations.